\begin{filecontents*}[overwrite]{cajas.sty}
% =====================================================================
%  cajas.sty -- Sistema semantico de cajas y convenciones TikZ
%  Extraido del Manual de Estilo LaTeX (USACH)
%  Compatible con article y beamer.
% =====================================================================
\NeedsTeXFormat{LaTeX2e}
\ProvidesPackage{cajas}[2026/07/31 Cajas semanticas y convenciones graficas]

% --- MATEMATICAS -----------------------------------------------------
\RequirePackage{amsmath,amssymb,amsthm,mathtools,bm}
\RequirePackage{cancel}

% --- REFERENCIAS -------------------------------------------------------
% [hyphens] permite que \url corte tras un guion -- sin esto, un slug
% largo sin "/" (tipo "camara-chilena-de-la-construccion-...") no tiene
% ningun punto de corte y se desborda del margen entero.
% En beamer, hyperref ya carga url.sty antes de llegar aqui; pedirlo de
% nuevo con otras opciones da "Option clash". Si ya esta cargado, no lo
% volvemos a pedir.
\makeatletter
\@ifpackageloaded{url}{}{\RequirePackage[hyphens]{url}}
\makeatother

% --- COLORES: table habilita \rowcolor -------------------------------
\RequirePackage[dvipsnames,table]{xcolor}

% --- CAJAS: [most] carga enhanced automaticamente --------------------
\RequirePackage[most]{tcolorbox}
\tcbuselibrary{breakable}

% --- GRAFICOS --------------------------------------------------------
\RequirePackage{tikz}
\usetikzlibrary{decorations.pathreplacing, arrows.meta}
\RequirePackage{pgfplots}
\pgfplotsset{compat=1.18}
\RequirePackage{forest}

% --- TABLAS Y LISTAS -------------------------------------------------
\RequirePackage{multirow}
\RequirePackage{booktabs}
\RequirePackage{array}
\RequirePackage{enumitem}
\RequirePackage{graphicx}

% =====================================================================
%  ESTILOS GLOBALES DE CAJAS
%  Mapa semantico:
%    definicion -> gris   -> conceptos y definiciones
%    teorema    -> verde  -> reglas, propiedades, tablas resumen
%    ejemplo    -> azul   -> ejemplos ilustrativos
%    nota       -> ambar  -> advertencias y conclusiones
%    enunciado  -> naranja-> comentes (verdadero/falso)
%    practica   -> teal   -> casos practicos numericos
%
%  POLITICA DE SALTO DE PAGINA (no breakable por defecto):
%  Ninguna caja es "breakable" por defecto. Si el contenido no cabe en lo
%  que resta de la pagina, la caja completa salta entera a la siguiente
%  en vez de partirse a la mitad (asi nunca queda un titulo huerfano ni
%  una lista cortada entre paginas).
%
%  Consecuencia: una caja MAS ALTA que una pagina no tiene adonde saltar
%  completa y se desborda. Por eso el contenido largo (una "Herramienta"
%  o "Receta" con varios pasos extensos, un ejemplo con 4+ Pasos de
%  desarrollo) NO debe ir dentro de una caja: sacalo del sistema de cajas
%  por completo (encabezado + texto corrido + enumerate normal). Ver la
%  seccion "Cuando NO usar una caja" en CLAUDE.md para el criterio exacto.
%
%  Si en un caso raro una caja SI necesita partirse (contenido que no
%  llega al umbral de "sacar de la caja" pero aun asi no cabe en una
%  pagina), agrega `breakable` a mano en esa instancia puntual:
%    \begin{tcolorbox}[teorema, breakable] ... \end{tcolorbox}
%  y verifica con pdftoppm que el corte cae entre parrafos, nunca a mitad
%  de un item de lista ni justo despues del titulo.
% =====================================================================
\tcbset{
  boxrule=0.6pt, arc=3pt,
  left=6pt, right=6pt, top=6pt, bottom=6pt,
  definicion/.style={
    enhanced,
    colback=gray!5!white,
    colframe=black!70,
    title=\textbf{Definición},
    fonttitle=\bfseries},
  teorema/.style={
    enhanced,
    colback=green!8!white,
    colframe=green!50!black,
    title=\textbf{Teorema},
    fonttitle=\bfseries},
  ejemplo/.style={
    enhanced,
    colback=blue!5!white,
    colframe=blue!70!black,
    title=\textbf{Ejemplo},
    fonttitle=\bfseries},
  nota/.style={
    enhanced,
    colback=yellow!10!white,
    colframe=orange!80!black,
    title=\textbf{Nota},
    fonttitle=\bfseries},
  enunciado/.style={
    enhanced,
    colback=orange!8!white,
    colframe=orange!85!black,
    fonttitle=\bfseries},
  practica/.style={
    enhanced,
    colback=teal!5!white,
    colframe=teal!60!black,
    fonttitle=\bfseries},
}

% =====================================================================
%  ATAJOS DE CONTENIDO
%  Encapsulan la estructura de un comente y de un caso practico
%  para que no haya que reescribir el andamiaje cada vez.
% =====================================================================

% \comente{titulo}{enunciado}{VERDADERO|FALSO}{desarrollo}
\newcommand{\comente}[4]{%
  \begin{tcolorbox}[enunciado, title={\large\textbf{Comentario:} #1}]
  \textbf{Enunciado:} \textit{``#2''}

  \medskip\hrule\medskip

  \textbf{Pauta de Respuesta --- \textcolor{red}{#3}}

  #4
  \end{tcolorbox}%
}

% \begin{caso}{titulo} ... \end{caso}
\newtcolorbox{caso}[1]{practica, title={\textbf{Caso Práctico:} #1}}

% \begin{conclusion} ... \end{conclusion}   (para anidar dentro de caso)
\newtcolorbox{conclusion}{nota, title=\textbf{Conclusión}}

% =====================================================================
%  CONVENCIONES DE COLOR PARA GRAFICOS
%  Estilos TikZ nombrados: en vez de recordar "blue!80!black, very thick"
%  se escribe [rp-original]. Cambiar el color aqui lo cambia en todo el curso.
% =====================================================================
\tikzset{
  eje/.style={->, thick},
  rp-original/.style={blue!80!black, very thick},
  rp-nueva/.style={red!80!black, very thick, dashed},
  rp-ficticia/.style={orange!80!black, dashed, very thick},
  ci-optima/.style={red!70!black, very thick},
  ci-otra/.style={black!50, thick},
  proyeccion/.style={dashed, gray, thick},
  senda/.style={green!50!black, dashed, thick},
  llave/.style={decorate, decoration={brace, amplitude=6pt}, thick},
  llave-inv/.style={decorate, decoration={brace, amplitude=6pt, mirror}, thick},
  paso/.style={draw, rounded corners, fill=orange!12!white,
               text width=5.5cm, align=center, inner sep=6pt,
               font=\footnotesize},
  flecha-paso/.style={->, very thick, orange!70!black},
  % Areas sombreadas frecuentes
  areaexcedente/.style={yellow!40, opacity=0.8},
  areaganancia/.style={blue!30, opacity=0.7},
  areanuevos/.style={green!40, opacity=0.8},
  areaperdida/.style={orange!40, opacity=0.8},
}

\endinput
\end{filecontents*}
\begin{filecontents*}[overwrite]{slides.sty}
% =====================================================================
%  slides.sty -- Estilo para presentaciones (documentclass beamer)
%  Reutiliza el mismo sistema de cajas que los apuntes.
% =====================================================================
\NeedsTeXFormat{LaTeX2e}
\ProvidesPackage{slides}[2026/07/31 Estilo de presentaciones]

\RequirePackage[spanish, es-noshorthands]{babel}
\RequirePackage[T1]{fontenc}
\RequirePackage{lmodern}

% cajas.sty NO carga geometry ni fancyhdr, por eso es seguro en beamer
\RequirePackage{cajas}

% Beamer trae su propio xcolor/tikz; estas lineas evitan choques
\setbeamertemplate{navigation symbols}{}
\setbeamertemplate{caption}[numbered]

% Las cajas tcolorbox dentro de beamer necesitan menos padding
\tcbset{
  left=4pt, right=4pt, top=4pt, bottom=4pt,
  fontupper=\small,
}

% Paleta del tema alineada a la semantica de cajas
\definecolor{acento}{RGB}{31,90,140}
\setbeamercolor{structure}{fg=acento}
\setbeamercolor{frametitle}{fg=acento}
\setbeamerfont{frametitle}{series=\bfseries}

% =====================================================================
%  Identidad y portada -- mismos comandos que estilo/apunte.sty, para
%  que un documento nuevo se complete igual sea apunte o presentacion.
% =====================================================================
\newcommand{\@curso}{Curso}
\newcommand{\curso}[1]{\renewcommand{\@curso}{#1}\subtitle{#1}}
\newcommand{\autoria}[1]{\author{#1}}
\newcommand{\institucion}[1]{\institute{#1}}

% \logo{figuras/logo.png.png} -- va solo en la portada (\titlegraphic),
% no en cada frame: repetirlo en cada slide compite con el frametitle.
% Si el archivo no existe el documento compila igual, como en apunte.sty.
\renewcommand{\logo}[1]{%
  \titlegraphic{\IfFileExists{#1}{\includegraphics[height=1.1cm]{#1}}{}}%
}

% --- RESALTADO DE PUNTOS CLAVE ---------------------------------------
% \destacar{texto} -- alias en espanol de \alert (respeta overlays,
% \destacar<2->{...} funciona igual que \alert<2->{...}). Tinta ambar,
% el mismo tono que la caja "nota" (orange!80!black en cajas.sty), para
% que resaltar en slides use la misma semantica de color que en apuntes.
\setbeamercolor{alerted text}{fg=orange!80!black}
\let\destacar\alert

% --- AGENDA RECURRENTE (opcional) ------------------------------------
% Llamar \agendaporsecciones en el preambulo de presentaciones largas
% (mas de ~15 slides o 3+ secciones). Inserta una slide de contenidos
% al empezar cada \section, resaltando la seccion en curso. En
% presentaciones cortas de una sola seccion no aporta nada: no activar.
\newcommand{\agendaporsecciones}{%
  \AtBeginSection[]{%
    \begin{frame}{Contenidos}
      \tableofcontents[currentsection]
    \end{frame}
  }%
}

\endinput
\end{filecontents*}
% =====================================================================
%  Econometría II -- Presentación Etapa 1 (jue 15-10-2026, máx. 10 min)
%  Fuente: informe-etapa-1-postulantes.tex (base de 108.251 x 30)
%  Lo que evalúa el profesor (clase 08-10): población y N, creación de
%  variables, y efecto esperado a priori de cada X sobre la Y.
% =====================================================================
\documentclass[aspectratio=169]{beamer}
\usetheme{default}
\usepackage{slides}

\autoria{David Cuevas \and Javier Li-Yau \and Benjamín Lizana \and Ana Nacif}
\institucion{Universidad de Santiago de Chile}

\title{Matrícula en las tres primeras preferencias}
\curso{Econometría II --- Trabajo práctico, Etapa 1 \quad Prof.\ Brian Pfeng}
\date{15 de octubre de 2026}

\begin{document}

\begin{frame}
  \titlepage
\end{frame}

% ---------------------------------------------------------------------
\section{Población}

\begin{frame}{¿De qué depende que un postulante se matricule en una de sus tres primeras preferencias?}
  \vfill
  \begin{center}
    {\Huge 55,7\,\%}\\[6pt]
    {\large de los 108.251 postulantes de la promoción 2025 lo logró}
  \end{center}
\end{frame}

\begin{frame}{Todos los datos vienen de cuatro archivos públicos del DEMRE, admisión 2026}
  \begin{center}
  \small
  \begin{tabular}{llr}
    \toprule
    \textbf{Archivo} & \textbf{Qué aporta} & \textbf{Filas} \\
    \midrule
    B: \texttt{ArchivoB\_Adm2026REG.csv} & Inscritos: colegio, sexo, ingreso, PACE, BEA & 317.700 \\
    C: \texttt{ArchivoC\_Adm2026REG.csv} & Notas, NEM, ranking y puntajes PAES & 317.700 \\
    D: \texttt{ArchivoD\_Adm2026REG.csv} & Postulaciones (una fila por carrera) & 1.686.689 \\
    M: \texttt{ArchivoMatr\_Adm2026.csv} & Matrícula: vía, tipo, preferencia, puntaje & 128.345 \\
    \bottomrule
  \end{tabular}
  \end{center}
  \medskip
  Se unen por \texttt{ID\_aux} en \destacar{una fila por estudiante}.
  \vfill
  {\footnotesize Fuente: DEMRE, Universidad de Chile, portal de bases de datos
  del proceso de admisión 2026, y su libro de códigos. Procesamiento propio
  en R (\texttt{tidyverse}, \texttt{writexl}).\par}
\end{frame}

\begin{frame}{Nuestra población son 108.251 egresados 2025 que postularon por la vía regular}
  \begin{center}
  \begin{tikzpicture}[x=1cm, y=0.62cm]
    \foreach \etiqueta/\ancho/\n [count=\i] in {
        {Unión de los cuatro archivos}/7.00/317.700,
        {Egresados en 2025}/4.82/218.989,
        {Alumnos regulares}/4.80/218.066,
        {Postularon por vía regular}/2.53/114.865,
        {No matriculados, o vía regular y tipo 1}/2.39/108.598} {
      \pgfmathsetmacro{\yy}{-\i}
      \fill[blue!25] (0,\yy-0.35) rectangle (\ancho,\yy+0.35);
      \node[left, font=\small] at (0,\yy) {\etiqueta};
      \node[right, font=\small] at (\ancho,\yy) {\n};
    }
    \fill[yellow!40] (0,-6.35) rectangle (2.39,-5.65);
    \node[left, font=\small\bfseries] at (0,-6) {Rindieron el proceso y NEM $>0$};
    \node[right, font=\small\bfseries] at (2.39,-6) {$N=108.251$};
  \end{tikzpicture}
  \end{center}
  \medskip
  \destacar{Los que no se matricularon se quedan}; salen los matriculados por
  otra vía o con otro tipo de matrícula.
\end{frame}

\begin{frame}{La $Y$ vale 0 si quedó en la 4.ª preferencia o más abajo, o si no se matriculó}
  \[
  \texttt{PREF\_TOP3\_DUMMY}_i =
  \begin{cases}
    1 & \text{se matriculó en su 1.ª, 2.ª o 3.ª preferencia} \\[2pt]
    0 & \text{quedó en la 4.ª o más abajo, o no se matriculó}
  \end{cases}
  \]
  \begin{center}
  \begin{tikzpicture}[x=1cm, y=1cm]
    \fill[green!45!black!35] (0,0) rectangle (6.13,1.0);
    \fill[orange!55]         (6.13,0) rectangle (6.98,1.0);
    \fill[red!35]            (6.98,0) rectangle (11.00,1.0);
    \node[align=center, font=\small] at (3.07,0.5) {$Y=1$: preferencia 1 a 3};
    \node[align=center, font=\small] at (8.99,0.5) {$Y=0$: no se matriculó};
    \draw[->] (6.55,1.0) -- (6.55,1.4) node[above, font=\small] {$Y=0$: preferencia 4 o más};
    \node[below, font=\small] at (3.07,0) {60.300 (55,7\,\%)};
    \node[below, font=\small] at (6.55,-0.4) {8.403 (7,8\,\%)};
    \node[below, font=\small] at (8.99,0) {39.548 (36,5\,\%)};
  \end{tikzpicture}
  \end{center}
\end{frame}

% ---------------------------------------------------------------------
\section{Variables creadas}

\begin{frame}{Del estudiante creamos tres dummies, además de la $Y$}
  \begin{center}
  \begin{tabular}{llr}
    \toprule
    \textbf{Variable} & \textbf{Construcción} & \textbf{Con 1} \\
    \midrule
    \texttt{TP\_DUMMY} & 1 si la rama es técnico profesional (T1 a T5) & 18.877 \\
    \texttt{PACE} & 1 si fue preseleccionado PACE (vacío $\to$ 0) & 7.790 \\
    \texttt{BEA} & 1 si tiene Beca de Excelencia Académica (vacío $\to$ 0) & 9.916 \\
    \bottomrule
  \end{tabular}
  \end{center}
  \medskip
  \begin{itemize}[itemsep=4pt]
    \item El ingreso familiar 99 («prefiero no responder») pasa a \texttt{NA}:
          29.199 casos (27\,\%).
    \item Un puntaje 0 en Lectora o Matemática 1 es quien no rindió esa prueba.
  \end{itemize}
\end{frame}

\begin{frame}{Del colegio creamos tres variables, calculadas por RBD antes de filtrar}
  \begin{center}
  \begin{tabular}{lll}
    \toprule
    \textbf{Variable} & \textbf{Cálculo por RBD} & \textbf{Mediana} \\
    \midrule
    \texttt{N\_ALUMNOS} & Inscritos del colegio en el proceso & 129 \\
    \texttt{INGRESO\_COLEGIO} & Mediana del tramo de ingreso (sin el 99) & 4 \\
    \texttt{N\_MATRICULADOS} & Matriculados vía regular, tipo 1 (0 si no hay) & 54 \\
    \bottomrule
  \end{tabular}
  \end{center}
  \medskip
  Se calculan con \destacar{todos los inscritos del colegio}, no solo con los
  que quedan en la base. Usamos la mediana porque los tramos de ingreso son
  ordinales y el último es abierto.
\end{frame}

% ---------------------------------------------------------------------
\section{Efecto esperado}

\begin{frame}{A mejor rendimiento académico, mayor probabilidad de quedar en las tres primeras}
  \begin{center}
  \begin{tabular}{lcl}
    \toprule
    \textbf{Variable} & \textbf{Signo} & \textbf{Por qué} \\
    \midrule
    \texttt{PTJE\_NEM}, \texttt{PTJE\_RANKING} & $+$ & Pesan en el puntaje ponderado \\
    \texttt{CLEC}, \texttt{MATE1} & $+$ & Más puntaje, más carreras alcanzables \\
    \texttt{BEA} (10\,\% mejor del curso) & $+$ & Buen rendimiento y beca para pagar \\
    \texttt{PACE} & $+$ & Cupos garantizados en universidades PACE \\
    \bottomrule
  \end{tabular}
  \end{center}
  \medskip
  El puntaje decide si \destacar{alcanza el corte} de las carreras que más
  quiere.
\end{frame}

\begin{frame}{El contexto socioeconómico debería subir la probabilidad de matricularse arriba}
  \begin{center}
  \begin{tabular}{lcl}
    \toprule
    \textbf{Variable} & \textbf{Signo} & \textbf{Por qué} \\
    \midrule
    \texttt{INGRESO\_PERCAPITA} & $+$ & Puede pagar y preparar la PAES \\
    \texttt{INGRESO\_COLEGIO} & $+$ & Entorno con más recursos y redes \\
    \texttt{GRUPO\_DEPENDENCIA} & $+$ & Particular pagado frente a público \\
    \texttt{TP\_DUMMY} & $-$ & Menos preparación PAES; más salida a CFT \\
    \bottomrule
  \end{tabular}
  \end{center}
  \medskip
  El efecto del ingreso puede ser menor por la \destacar{gratuidad}: es
  parte de lo que el modelo tiene que medir.
\end{frame}

\begin{frame}{En el colegio y el territorio, el signo no es obvio a priori}
  \begin{center}
  \begin{tabular}{lcl}
    \toprule
    \textbf{Variable} & \textbf{Signo} & \textbf{Por qué} \\
    \midrule
    \texttt{N\_MATRICULADOS} & $+$ & Cultura de seguir estudios en el colegio \\
    \texttt{N\_ALUMNOS} & ? & Más recursos, pero menos atención por alumno \\
    \texttt{CODIGO\_REGION} & ? & Más oferta en Santiago, pero más competencia \\
    \texttt{SEXO} & ? & Mujeres: más NEM; hombres: más MATE1 \\
    \bottomrule
  \end{tabular}
  \end{center}
  \medskip
  \texttt{N\_MATRICULADOS} se lee mejor como tasa: \destacar{matriculados
  sobre inscritos} del colegio.
\end{frame}

\begin{frame}{Las variables de la matrícula no pueden ser $X$: solo existen si se matriculó}
  \begin{center}
  \begin{tabular}{ll}
    \toprule
    \textbf{Variable} & \textbf{Problema} \\
    \midrule
    \texttt{PREFERENCIA} & Con ella se construye la $Y$ \\
    \texttt{PTJE\_POND} & Es el puntaje de la carrera en que ya quedó \\
    \texttt{VIA}, \texttt{TIPO\_MATRICULA} & Vacías en los 39.548 no matriculados \\
    \bottomrule
  \end{tabular}
  \end{center}
  \medskip
  Usarlas sería explicar la $Y$ \destacar{con la misma $Y$}: quedan en la
  base solo como información.
\end{frame}

% ---------------------------------------------------------------------
\begin{frame}{Lo que deja la Etapa 1 y lo que viene}
  \begin{itemize}[itemsep=7pt]
    \item Población: $N = 108.251$ egresados 2025 que postularon por la vía
          regular, incluidos los no matriculados.
    \item La $Y$ queda en 55,7\,\% frente a 44,3\,\%.
    \item El rendimiento y el ingreso deberían subir la probabilidad; TP la
          bajaría; colegio y región quedan abiertos.
  \end{itemize}
  \vfill
  \begin{tcolorbox}[nota]
  \textbf{Etapa 2:} cruzar cada $X$ con la $Y$ y agregar variables del
  MINEDUC por RBD (cursos, alumnos por profesor, ruralidad).
  \end{tcolorbox}
\end{frame}

% =====================================================================
\appendix
\section*{Respaldo}

\begin{frame}{Respaldo: estadísticas de la base final}
  \begin{center}
  \begin{tabular}{lrrrr}
    \toprule
    \textbf{Variable} & \textbf{Mín.} & \textbf{Mediana} & \textbf{Promedio} & \textbf{Máx.} \\
    \midrule
    \texttt{PTJE\_NEM}          & 163 & 779 & 767 & 1.000 \\
    \texttt{PTJE\_RANKING}      & 163 & 793 & 781 & 1.000 \\
    \texttt{CLEC\_REG\_ACTUAL}  & 0   & 634 & 640 & 1.000 \\
    \texttt{MATE1\_REG\_ACTUAL} & 0   & 645 & 668 & 1.000 \\
    \texttt{INGRESO\_COLEGIO}   & 1,0 & 4,0 & 4,8 & 10,0 \\
    \bottomrule
  \end{tabular}
  \end{center}
  \medskip
  3.196 colegios; 59.799 mujeres (55,2\,\%) y 48.452 hombres.
\end{frame}

\begin{frame}{Respaldo: resto de las variables continuas}
  \begin{center}
  \begin{tabular}{lrrrr}
    \toprule
    \textbf{Variable} & \textbf{Mín.} & \textbf{Mediana} & \textbf{Promedio} & \textbf{Máx.} \\
    \midrule
    \texttt{PROMEDIO\_NOTAS}           & 4,20  & 6,23  & 6,19  & 7,00 \\
    \texttt{PTJE\_POND} (matriculados) & 354,6 & 710,7 & 713,6 & 998,4 \\
    \texttt{N\_ALUMNOS}                & 8     & 129   & 164   & 947 \\
    \texttt{N\_MATRICULADOS}           & 0     & 54    & 70    & 447 \\
    \bottomrule
  \end{tabular}
  \end{center}
\end{frame}

\end{document}
